% =============================================================
%  front/abstract.tex  --  Abstract / Kurzfassung
% =============================================================

\chapter*{Über dieses Dokument}
\addcontentsline{toc}{chapter}{Über dieses Dokument}

\begin{infobox}[title=Wie dieses Dokument zu lesen ist]
Das Handbuch ist in sechs Teile gegliedert.
\textbf{Teil~I} legt die Theorie fest (ROS\,2, Kalman/EKF, Odometrie, Stereosehen,
Kinematik, Dynamik, Trajektorien, Schrittmotoren, Zykloidgetriebe).
\textbf{Teil~II} behandelt die Hardware-Auswahl inklusive Strom-/Spannungsmessung
zur Drehmomentschätzung und PCB-Optimierung.
\textbf{Teil~III} ist die Schritt-für-Schritt-Installationsanleitung (VM, Linux,
ROS\,2).
\textbf{Teil~IV} erklärt die Simulation, den CAD-zu-URDF-Workflow und die
entscheidende Frage, wo sich Simulation und reale Implementierung trennen.
\textbf{Teil~V} arbeitet die vier Projekte mit Knoten-Architektur und
\code{rqt\_graph} durch.
\textbf{Teil~VI} ist die praktische Stereokamera-Integration mit Codebeispielen.
Den Abschluss bildet eine empfohlene Dokumentationsmethodik.
\end{infobox}
